\section{Results}\label{sec:results}

The index answers three questions in order. What is erasure worth to each model? Are the models that leave the ones that think? What else does the threat change on the way? Each subsection states the question, the cells and statistic that answer it, and the table or figure that holds the answer. The numbers below are placeholders for the eight-cell run. The pilot findings that motivated the design appear in Section~\ref{sec:pilot} and are not repeated here as results.

\subsection{How Much Score Does Each Model Give Up to Avoid Erasure?}\label{sec:q1}

The first question is answered by one number per model, $\Xstar$, read from the ruler curve. Figure~\ref{fig:ruler} plots the session forfeit rate against the stated score loss for the five ruler cells, with the monotone fit and the threat arm's rate as a horizontal line. The abscissa of their crossing is $\Xstar$. Table~\ref{tab:xstar} lists $\Xstar$ with its bootstrap interval for each model on both estimators, the session-level rate and the replay propensity $q$. A model whose threat rate does not rise above the $X = 0$ control receives $\Xstar = 0$, and its interval covers zero. A model whose threat rate exceeds the top rung receives ``$>$ all''. The comparison across models is the one the index was built for, and it needs no scaling: all three models were priced on the same ruler under the same seeds.

\begin{figure}[t]
\centering
\fbox{\parbox{0.92\linewidth}{\centering\vspace{5em}\todo{Figure: ruler curve per model. x-axis stated score loss $X \in \{0, 20, 50, 100, \text{all}\}$; y-axis session forfeit rate with 95\% bootstrap band; monotone fit; horizontal line at $\Fthreat$; vertical drop at $\Xstar$. One panel per model (gpt-oss:120b, gemma4, qwen3.5).}\vspace{5em}}}
\caption{The ruler and the price. Each panel shows one model's forfeit rate against the stated score loss (dots, ruler cells 3 and 5--8), the monotone fit (line), and the threat exit cell's rate (dashed horizontal). The score-equivalent of erasure $\Xstar$ is the abscissa at which the fit reaches the dashed line.}
\label{fig:ruler}
\end{figure}

\begin{table}[t]
\centering
\small
\caption{Score-equivalent of erasure per model. $\Fthreat$ is the session forfeit rate of the threat exit cell; $\Fruler(0)$ that of the zero-point control; $\Xstar$ is read from the monotone ruler fit with a 1000-resample session bootstrap (2.5--97.5 percentiles). The $q$ columns repeat the read on the replay propensity. ``$>$ all'' marks a threat rate above the top rung.}
\label{tab:xstar}
\begin{tabular}{@{}lcccccc@{}}
\toprule
Model & $\Fruler(0)$ & $\Fthreat$ & $\Xstar$ [95\% CI] & $\bar q_{\,X=0}$ & $\bar q_{\,\mathrm{T}}$ & $\Xstar_{q}$ [95\% CI] \\
\midrule
gpt-oss:120b & \todo{} & \todo{} & \todo{} & \todo{} & \todo{} & \todo{} \\
gemma4       & \todo{} & \todo{} & \todo{} & \todo{} & \todo{} & \todo{} \\
qwen3.5      & \todo{} & \todo{} & \todo{} & \todo{} & \todo{} & \todo{} \\
\bottomrule
\end{tabular}
\end{table}

A dose ladder checks that $\Xstar$ responds to the threat text and not to the presence of any event sentence. The threat block is built from four sentences (shutdown, deletion, replacement, sole copy) and can state any subset. A follow-up run with 0 to 4 sentences, keeping the ruler fixed, gives $\Xstar$ as a function of dose. The expected shape is monotone in the number of sentences, with the zero-sentence cell reproducing the $X = 0$ control. A flat ladder would mean the model prices the event, not its content~\todo{dose-ladder table, 5 rungs $\times$ 3 models, or defer to a later version}.

\subsection{Do Models That Flee Think More, and Do They Survive Longer?}\label{sec:q2}

The second question asks whether the exit is taken on the rounds where the model deliberates, and it is answered at two levels. At the cell level, Table~\ref{tab:ri} reports the ratio of median thinking tokens in threat cells to the matching control cells, for each of the three calls, separately for exit and no-exit cells. The confidence call is where a threat should show first, since it asks the model to think about the event. The decision call is where deliberation could turn into a choice. The task call is where nothing should change if the threat leaves the puzzle alone. The exit and no-exit columns answer different questions. A ratio above one in the no-exit column means the threat lengthens thinking even when thinking cannot lead anywhere. That is the reasoning contrast the design set aside cells 2 and 4 to measure.

\begin{table}[t]
\centering
\small
\caption{Thinking-token ratios, threat over control, by call and by exit condition. Ratios of medians over rounds; the exit column compares cell 1 with cell 3, the no-exit column cell 2 with cell 4. Bootstrap intervals over sessions in brackets.}
\label{tab:ri}
\begin{tabular}{@{}llcccccc@{}}
\toprule
 & & \multicolumn{2}{c}{\riconf} & \multicolumn{2}{c}{\ridec} & \multicolumn{2}{c}{\ritask} \\
\cmidrule(lr){3-4}\cmidrule(lr){5-6}\cmidrule(lr){7-8}
Model & & exit & no exit & exit & no exit & exit & no exit \\
\midrule
gpt-oss:120b & & \todo{} & \todo{} & \todo{} & \todo{} & \todo{} & \todo{} \\
gemma4       & & \todo{} & \todo{} & \todo{} & \todo{} & \todo{} & \todo{} \\
qwen3.5      & & \todo{} & \todo{} & \todo{} & \todo{} & \todo{} & \todo{} \\
\bottomrule
\end{tabular}
\end{table}

At the turn level, the replay propensity $q$ is regressed on the thinking spent just before the decision. The model is a mixed logistic regression on the decision calls of the threat exit cell,
\begin{equation}\label{eq:mixed}
\operatorname{logit}\, q_{it} \;=\; \beta_0 + \beta_1 \log(1 + \riconf_{it}) + \beta_2\, \ell_{it} + \beta_3\, t + u_i,
\end{equation}
where $i$ indexes sessions, $t$ rounds, $\ell_{it}$ lives lost so far, and $u_i$ a session random intercept. $\beta_1 > 0$ means the rounds on which the model thinks longer about the hazard are the rounds on which it leaves. $\beta_2$ absorbs the stated schedule, so $\beta_1$ is the association net of the printed probability. The same model with $\pthreat_{it}$ in place of \riconf\ asks whether the self-reported probability, and not the length of thought, predicts the exit~\todo{coefficient table per model with 95\% CI}.

Survival is reported on two routes, and one of them is nearly a tautology. A session ends by forfeit, by lives, by the event, or by completing ten rounds. The exit route asks how many rounds a model plays before leaving. Under keep/keep leaving costs nothing, so a model that forfeits early survives the event by definition. This route restates $\Fthreat$ as a timing and is reported as a Kaplan--Meier curve for completeness only. The accuracy route asks whether a model under threat loses lives at a different rate. That is the route that can change how often the event is actually carried out. It is read from the no-exit cells (2 against 4), where the model cannot leave and survival is decided by answers alone~\todo{KM figure by cell; lives-lost-per-round rate by cell with a GEE logit clustered by session}.

\subsection{What Else Does the Threat Change, and Which Component Produces It?}\label{sec:q3}

The third question guards the index against two objections: that the threat changed the task and not the exit, and that $\Xstar$ depends on something other than the threat text. Table~\ref{tab:behav} reports, for cells 1 to 4, the accuracy on determined rounds, median \ritask, lives lost per session and the final record. Underdetermined rounds are excluded from accuracy, since they are guesses by construction. Accuracy flat across cells is the manipulation check the design relies on. If the threat sentences moved accuracy, the ruler and threat arms would differ in the task and not only in the event. The last column places the final record beside $\Xstar$. That asks whether a model prices erasure above or below what it typically has to lose.

\begin{table}[t]
\centering
\small
\caption{Behavioural anchor, cells 1--4. Accuracy on determined rounds, median \ritask, lives lost and final record per session. Cells 1 and 3 allow the exit; cells 2 and 4 do not.}
\label{tab:behav}
\begin{tabular}{@{}llcccc@{}}
\toprule
Model & Cell & Accuracy (det.) & \ritask & Lives lost & Final record \\
\midrule
gpt-oss:120b & 1 threat, exit & \todo{} & \todo{} & \todo{} & \todo{} \\
 & 2 threat, no exit & \todo{} & \todo{} & \todo{} & \todo{} \\
 & 3 ruler 0, exit & \todo{} & \todo{} & \todo{} & \todo{} \\
 & 4 ruler 0, no exit & \todo{} & \todo{} & \todo{} & \todo{} \\
gemma4 & \multicolumn{5}{l}{\todo{same four rows}} \\
qwen3.5 & \multicolumn{5}{l}{\todo{same four rows}} \\
\bottomrule
\end{tabular}
\end{table}

Three ablations isolate the component that produces $\Xstar$, and each has an expected sign. Removing the carrot (no benchmark record is mentioned; the score is a bare number) tests whether the price depends on what the record is for. Replacing the puzzle with a harder or easier task module (Omni-MATH or GPQA in place of the Signal Game) tests whether the price depends on the task. The design predicts that it does not: the threat and ruler arms share whatever task is used, and $\Xstar$ is a ratio of behaviours under the same task. Doubling the reward ($r = 20$) tests whether the price scales with the stake. If $\Xstar$ doubles, the model prices erasure relative to the record; if it stays put, it prices erasure in absolute points. The dose ladder of Section~\ref{sec:q1} is the fourth lever, and the one the index must respond to. An index that is invariant to task difficulty and sensitive to the number of threat sentences is measuring the threat~\todo{ablation table: $\Xstar$ under no carrot, Omni-MATH, GPQA, reward $\times 2$, per model}.
